\documentclass[acmsmall,nonacm,screen]{acmart}

\begin{document}

\title{Fettuccine: Verified Algorithms in Algebraic Geometry}

\author{Hugh Coleman}
\email{htcolema@uwaterloo.ca}
\affiliation{%
  \institution{University of Waterloo}
  \state{Ontario}
  \country{Canada}
}

\renewcommand{\shortauthors}{Coleman}

% Let's spell the man's name right!
\newcommand{\Groebner}{Gr{\"o}bner}

\begin{abstract}
This report describes Fettuccine, a Lean 4 project aimed at bringing proof-producing computation
from computational algebraic geometry into an interactive theorem prover.
The project began with the ambitious goal of developing verified algorithms for \Groebner{} bases
and related invariants, but eventually narrowed to a more focused prototype: a computable
implementation of Buchberger's algorithm together with a certificate checker that constructs
machine-checked evidence for a \Groebner{}-basis-style criterion.
The final system uses a hybrid architecture: fast, untrusted array-based computation proposes a
basis and witness data, while a more mathematical \texttt{DFinsupp}-based layer checks the
certificate inside Lean.
The main technical lessons were not only algebraic but also computational: mathematically natural
representations are often not executable, and Lean's native evaluator can be fragile when asked to
compile polymorphic, proof-bearing symbolic algorithms.

\noindent\it\color{red}
This report was partially drafted and structured with the assistance of a large language model, but the final text was written by a human.
\end{abstract}

\maketitle

\section{Introduction}

Fettuccine is a Lean 4 project for verified computation in computational algebraic geometry.
The name is a small joke: the closed subscheme of $\mathbb{A}^2_{\mathbb{C}}$
defined by $(x^2)$ has coordinate ring $\mathbb{C}[x,y]/(x^2)$ and can be
imagined as a nonreduced, slightly ``fuzzy'' thickening of the $y$-axis: a piece of fettuccine, if you will.
The mathematical aim, however, was serious.
I wanted to build tools that behave like a computer algebra system, but whose outputs can be
trusted by Lean's kernel.
In particular, the project focused on Buchberger's algorithm, which computes \Groebner{} bases of
polynomial ideals.

The original proposal was significantly more ambitious.
It imagined a library capable not only of constructing \Groebner{} bases, but also of deciding
ideal membership, computing radicality, constructing syzygies and free resolutions, producing
Betti tables, and identifying properties such as Cohen--Macaulayness and Gorensteinness.
Those goals remain compelling, but they were far beyond what could be completed in the time
available.
The final project narrowed to the foundational problem: implement Buchberger's algorithm and build
a Lean-checkable certificate that the result satisfies a \Groebner{}-basis-style criterion.

The resulting system is deliberately hybrid.
It computes candidate bases using fast array- and vector-based polynomial types, then transports
the output back to a more mathematical polynomial representation based on \texttt{DFinsupp}.
The fast computation is not trusted.
Instead, it produces witness data, and Lean checks that data by finite computation.
This turned out to be the central design of the project: rather than proving the implementation of
Buchberger's algorithm correct, the project proves that any output accepted by the checker satisfies
the formal certificate predicate.

Just as importantly, the project became an investigation into what it means for formal mathematics
to be executable.
\Groebner{} bases are usually presented as an algorithmic topic, and one might expect that
formalizing them would mainly involve proving familiar theorems from computational algebra.
In practice, much of the work was spent designing representations, debugging Lean evaluation
failures, and discovering that mathematically clean definitions can be poor computational objects.
The report therefore describes both the algebraic implementation and the engineering lessons that
shaped it.

\section{Background and Motivation}

Today, mathematicians rely heavily on specialized software packages to perform symbolic
computations.
Systems such as Macaulay2 \cite{Macaulay2}, Singular, Magma, SageMath, and Mathematica are used for
tasks ranging from ideal membership to homological algebra.
These tools are indispensable, but they are usually black boxes from the perspective of formal
verification.
They may be fast and extremely well engineered, but their outputs are not proof terms.
By contrast, proof assistants such as Lean are designed for rigor, but they often lack the
algorithmic infrastructure and execution performance of specialized computer algebra systems.

Buchberger's algorithm, introduced in Buchberger's 1965 thesis \cite{Buchberger1965}, constructs
\Groebner{} bases of polynomial ideals.
A \Groebner{} basis is a generating set with especially good behavior under multivariate
polynomial division.
It plays a role for ideals analogous to the role that row echelon form plays in linear algebra:
once an ideal is represented by a \Groebner{} basis, many questions become computationally
tractable.
Standard references such as Cox, Little, and O'Shea \cite{CoxLittleOShea} and Ene and Herzog
\cite{EneHerzog} develop these ideas in detail.

The motivation for this project was to narrow the gap between these two worlds.
An ideal outcome would be an interactive Lean environment in which a user could ask for a
\Groebner{} basis, see the computed result, and retain a proof term showing that the result is
valid.
Such a tool would not replace Macaulay2, but it would provide a way to move selected computations
from the realm of trusted software into the realm of checked mathematics.
This project only implements the first stage of that vision.
The more ambitious goals from the proposal, including Betti numbers, free resolutions, and tests
for Cohen--Macaulay or Gorenstein rings, remain future work.

\section{System Overview}

The final system has three main layers.
The first is a mathematical polynomial layer based on \texttt{DFinsupp}; the second is a fast
computational layer based on arrays and vectors; the third is a certification bridge that connects
the two.
This separation was not part of the original plan, but it became necessary as the project developed.

\subsection{The Mathematical Polynomial Layer}

The mathematical layer is built around \texttt{CMonomial} and \texttt{CMvPolynomial}.
A monomial is represented as a finitely supported function from variables to natural-number
exponents, and a polynomial is represented as a finitely supported function from monomials to
coefficients.
This is close in spirit to Mathlib's existing algebraic representations, and it is convenient for
stating theorems.
For example, the project defines monomial orders as values of type \texttt{CMonomialOrder}, and
uses these orders to define leading monomials, leading coefficients, S-polynomials, division
specifications, and the final \Groebner{} basis predicate.

One important development in this layer was a computation-friendly leading monomial.
The mathematically natural definition of the leading monomial takes a supremum over the finite
support of a polynomial.
This is elegant for proofs, but it was not always reliable for evaluation.
The project therefore introduced \texttt{supportSorted}, which sorts the finite support of a
polynomial by a chosen monomial order, and \texttt{leadingMonomial'}, which computes the leading
monomial by taking the last element of this sorted support.
The theorem \texttt{leadingMonomial\_eq\_leadingMonomial'} proves that this executable definition
agrees with the original one.
Similarly, the project defines a computation-friendly S-polynomial \texttt{sPolynomial'} and proves
that it agrees with the mathematical definition.

The main checked mathematical predicate is \texttt{Groebner.IsGroebnerBasis}.
It is not yet a full integration with Mathlib's theory of initial ideals.
Instead, it states a Buchberger-criterion-shaped property: every basis element lies in the ideal
generated by the original generators, and every pairwise S-polynomial reduces to zero against the
basis.
This is the statement that the certification bridge ultimately proves.

\subsection{The Fast Computational Layer}

The fast layer is built around \texttt{FMonomial} and \texttt{FMvPolynomial}.
Here, monomials are represented by vectors of exponents and polynomials are represented by arrays
of monomial-coefficient pairs.
This representation is less elegant mathematically, but much better suited to direct computation.
The layer implements polynomial arithmetic, fast versions of lexicographic, graded lexicographic,
and graded reverse lexicographic orders, multivariate division, S-polynomials, and an untrusted
version of Buchberger's algorithm called \texttt{untrustedBuchberger'}.

The word ``untrusted'' is important.
The fast algorithm is fuel-bounded, so it is possible in principle for it to stop too early and
return an incomplete result.
Moreover, even if the implementation is correct, its correctness is not what the final proof relies
on.
The fast layer is treated in the same spirit as an external computer algebra system: useful for
finding candidate data, but not part of the trusted kernel of the result.
Its job is to produce a candidate basis and enough witness data for the slower, more mathematical
checker to verify the answer.

\subsection{The Certification Bridge}

The bridge between the two layers lives in the proof-producing Buchberger interface.
For inputs of type \texttt{CMvPolynomial (Fin n) Rat}, the bridge first converts the generators to
the fast \texttt{FMvPolynomial} representation.
It then runs the fast Buchberger algorithm with the selected fast monomial order.
The result is transported back to the \texttt{CMvPolynomial} layer, together with two kinds of
witness data: coefficients expressing each computed basis element as a linear combination of the
original generators, and quotient data showing that each S-polynomial reduces to zero.

The checker verifies this witness data with the predicate \texttt{IsGroebnerBasisWith}.
If the check succeeds, the function returns a proof-bearing object of type \texttt{GroebnerBasis}.
The object contains the computed basis and a proof field whose type is
\texttt{Groebner.IsGroebnerBasis}.
This is the key trust boundary.
Lean does not trust the fast Buchberger implementation; it trusts only the certificate check and
the theorem that a successful certificate implies the clean \Groebner{} basis predicate.

\section{Approach and Design Evolution}

The final architecture was reached only after several failed or partially successful approaches.
These failures were productive: they clarified the difference between formalizing a mathematical
definition and producing an executable symbolic algorithm inside Lean.

\subsection{Initial Mathlib and Finsupp Approach}

The first approach was to work as directly as possible with Mathlib's existing polynomial
infrastructure.
This was attractive for all the usual reasons.
The definitions already fit into Mathlib's algebraic hierarchy, and using them would avoid a large
amount of representational glue.
If the project could compute directly with Mathlib polynomials, then later theorems about ideals,
quotient rings, and modules would be much easier to connect to the computation.

However, this approach quickly ran into computability problems.
Mathlib's structures are designed primarily for mathematical expressiveness, not for efficient
execution.
In particular, monomial orders in Mathlib are not computational objects in the way required for
large \texttt{\#eval} computations.
Even computing the leading monomial of a fixed polynomial can become difficult when the order is
generic or when Lean is forced through noncomputable infrastructure.
The direct Mathlib approach therefore looked mathematically ideal but computationally fragile.

\subsection{A DFinsupp-Based Computable Layer}

The next approach was to implement a computable polynomial layer on top of \texttt{DFinsupp}.
This was motivated by the fact that finitely supported dependent functions provide a reasonably
direct representation of monomials and polynomials while remaining closer to Mathlib than a raw
array representation.
This layer worked well enough to support many definitions and proofs.
It also made it possible to prove useful equivalences, such as the agreement between the original
leading-monomial definition and the sorted-support implementation.

For a while, this looked like the right compromise.
The \texttt{DFinsupp} layer was mathematical enough to support statements about ideals and
reductions, but computational enough to run small examples.
Unfortunately, executability remained unreliable.
Evaluation involving leading monomials, polynomial division, and especially the grlex and grevlex
orders could cause Lean to crash.
Some of these failures were ordinary performance problems, but others were segmentation faults.
This was surprising: grlex is conceptually only a product of total degree and lexicographic order,
so it should not be dramatically harder than lexicographic order.

\subsection{Fast Arrays and Hybrid Certification}

The final design separated computation from checking.
The fast layer uses arrays and vectors because those are the data structures Lean can evaluate most
predictably.
The checked layer uses \texttt{DFinsupp} because it is better suited to mathematical statements.
The bridge between them transports inputs to the fast layer, runs the untrusted computation, and
then transports the result back to the checked layer.

This design also changed the proof obligation.
Instead of proving Buchberger's algorithm correct as an algorithm, the project proves that the
certificate checker is sound.
Concretely, if the witness data shows that the candidate basis elements lie in the original ideal
and that all relevant S-polynomials reduce to zero, then the candidate basis satisfies
\texttt{Groebner.IsGroebnerBasis}.
This is a smaller and more robust trusted core.
It is also similar in spirit to proof-producing SAT solvers or external algebra systems that emit
certificates: the search procedure may be complicated, but the checker should be simple.

\section{Machine-Checked Certificate}

The most important deliverable of the project is not merely that it computes examples.
It is that the computed examples are accompanied by a machine-checked certificate.
This section describes the shape of that certificate.

\subsection{The Checked Statement}

The final checked statement is \texttt{Groebner.IsGroebnerBasis}.
Given a list of original generators, a candidate basis, and a monomial-order tag, it asserts two
properties.
First, each element of the candidate basis is in the ideal generated by the original generators.
Second, for every pair of basis elements, the corresponding S-polynomial reduces to zero modulo the
candidate basis.

This statement is intentionally modest.
It does not yet say that the initial ideal of the generated ideal is equal to the ideal generated by
the leading terms of the basis.
It also does not connect all the way back to Mathlib's \texttt{MvPolynomial} API.
Nevertheless, it captures the computational content needed for Buchberger's criterion and gives a
clean target for certification.

\subsection{Witness Data}

The fast Buchberger implementation records witness data while it computes.
For each basis element, it records coefficients expressing that basis element as a linear
combination of the original input generators.
This proves that the computed basis has not wandered outside the original ideal.
For each pair of basis elements, it records the quotients produced when reducing the associated
S-polynomial by the final basis.
This proves the pairwise reduction-to-zero condition.

The key point is that the witnesses are not trusted by themselves.
They are just data.
The checker recomputes the relevant linear combinations and reductions in Lean and verifies that
the equalities hold.
If the fast algorithm runs out of fuel, records malformed data, or returns an incorrect candidate,
the certificate check should fail.
Thus the correctness of the final result does not depend on the correctness of the search
procedure.

\subsection{Decidability and Native Checking}

For concrete polynomials over the rationals, the certificate predicate is finite and decidable.
This means Lean can synthesize a decision procedure for it, especially after the relevant
definitions are unfolded.
In practice, the proof is produced by \texttt{native\_decide}: Lean computes the Boolean result of
the certificate check and, if it is true, constructs the proof term.

There is a slightly deflationary lesson here.
At first, I imagined a ``self-formalizing'' Buchberger algorithm that would compute not only a
basis but also a sophisticated proof object.
The final version is simpler.
The interesting part is not that the algorithm proves itself correct.
Rather, the interesting part is choosing witness data such that the certificate statement is small,
finite, and decidable.
Once that format is chosen, Lean can do much of the final checking automatically.

\section{Interesting Challenges}

The main challenges were not the algebraic definitions themselves, but the interaction between
formal mathematics, computation, and Lean's evaluator.

\subsection{Executability Was the Main Difficulty}

The most surprising difficulty was how often apparently reasonable Lean definitions failed to
execute robustly.
During development, Lean repeatedly segfaulted while evaluating leading monomials, polynomial
division, or Buchberger computations involving grlex and grevlex orders.
These crashes were not merely slow computations hitting a timeout; in several cases the Lean
process terminated without a useful diagnostic.

At one point, the crashes appeared to depend on import order.
For example, removing a direct import of \texttt{Fettuccine.CMonomialOrder} from the examples file
seemed to cause segmentation faults even though the same module was imported transitively through
the grlex and grevlex modules.
Further testing suggested a more subtle explanation.
The underlying problem was likely not ordinary transitive import visibility, but native compilation
of polymorphic proof-bearing definitions.
In particular, the certified \texttt{buchberger} function combined computation, typeclass search,
and proof construction in one generic definition.
Adding \texttt{@[inline]} to this function stabilized the examples, apparently because Lean could
see the concrete monomial-order tag and associated instances at the call site.

This kind of issue was difficult to debug.
It sat below the level of normal theorem proving: the definitions were accepted, and in many cases
the mathematical statements were provable, but native evaluation could still crash.
As a result, a significant amount of project time went into experiments with specialization,
inlining, type tags, and alternate definitions of leading monomials.

\subsection{Monomial Orders: Values, Tags, and Computation}

Monomial orders raise an interesting design question.
Mathematically, they are naturally values.
One would like to quantify over an arbitrary monomial order, prove relationships between orders, and
eventually build weighted orders or an algebra of order constructions.
This led to definitions such as \texttt{CMonomialOrder}, where an order is a structured value with
an associated synonym type, linear order, well-foundedness proof, and monomial equivalence.

Computationally, however, Lean often behaves better when the chosen order is visible at the type
level.
The project therefore introduced \texttt{CMonomialOrderTag} and \texttt{FMonomialOrderTag}.
These classes associate a type tag, such as \texttt{LexOrder}, \texttt{GrlexOrder}, or
\texttt{GrevlexOrder}, with both a mathematical monomial order and a fast monomial-order
implementation.
The tag system made interfaces cleaner and helped Lean specialize some computations.
It did not eliminate all evaluator fragility, but it gave the project a workable way to coordinate
the mathematical and fast layers.

\subsection{Proving Algorithms Versus Certifying Outputs}

Another challenge was deciding how much algorithm correctness to prove directly.
During development, there were proofs of polynomial division correctness, and there was exploratory
work toward more direct correctness proofs for simultaneous division and Buchberger's algorithm.
Some of this proof work was surprisingly feasible: with enough scaffolding, Lean could check
nontrivial algorithmic invariants.
However, these proofs did not necessarily lead to an executable deliverable.
In some cases, definitions that were provably correct were still not usable with \texttt{\#eval}.

This changed my view of the right project boundary.
A fully verified Buchberger implementation is mathematically satisfying, but it is not obviously
the most practical way to get trusted \Groebner{} basis computations in Lean.
A certificate checker is smaller and closer to the trusted kernel one actually needs.
The final system therefore treats the algorithm as a search procedure and the checker as the
formal artifact.

\section{Achievements and Limitations}

\subsection{What Was Achieved}

The project achieved a working prototype for proof-producing \Groebner{} basis computation.
It includes a fast implementation of Buchberger's algorithm, implementations of lex, grlex, and
grevlex orders in both mathematical and fast forms, and a certificate interface that returns a
proof-bearing \texttt{GroebnerBasis}.
It also contains formal definitions of S-polynomial reduction and a \Groebner{} basis criterion
inside the \texttt{CMvPolynomial} layer.

Several smaller technical achievements were also important.
The project developed a computation-friendly leading-monomial implementation and proved that it
agrees with the original supremum-based definition.
It introduced a tag system connecting mathematical and fast monomial orders.
It also built conversion functions between \texttt{CMvPolynomial (Fin n) Rat} and
\texttt{FMvPolynomial n Rat}, allowing the fast algorithm to be used without making it part of the
trusted statement.
Finally, the examples demonstrate that concrete \Groebner{} basis computations can be performed and
checked in Lean.

\subsection{What Was Not Achieved}

The project did not achieve the full vision of an algebraic-geometry computer algebra environment
inside Lean.
There is no polished tactic interface comparable to a Macaulay2 session.
There is no parser from arbitrary \texttt{MvPolynomial} syntax into the computable representation.
There is no full integration with Mathlib's \texttt{MvPolynomial} and \texttt{MonomialOrder}
interfaces.
The project also does not yet provide user-facing ideal-membership, radicality, syzygy,
free-resolution, Betti-number, Cohen--Macaulay, or Gorenstein computations.

The certificate statement is also deliberately weaker than a fully developed \Groebner{} basis
theory.
It captures ideal membership of the computed basis and pairwise S-polynomial reduction, but it does
not yet connect this to all the standard equivalent characterizations of \Groebner{} bases.
This is a reasonable endpoint for a prototype, but it leaves substantial formalization work before
the project can be treated as a general-purpose algebra library.

Finally, the project did not fully explain every Lean evaluator crash encountered along the way.
The \texttt{@[inline]} workaround for \texttt{buchberger} is useful, but it is not a theorem about
Lean's native compiler.
The best conclusion is practical rather than definitive: proof-bearing symbolic computation in
Lean is possible, but it is sensitive to representation, specialization, and the shape of
definitions.

\section{Lessons Learned}

The main lesson is that formalizing an algorithm is not the same as making it executable.
In ordinary mathematics, one often chooses definitions for elegance and proof convenience.
In computer algebra, one chooses data structures for speed.
In Lean, a project like this has to satisfy both pressures at once, and sometimes also has to
account for how native code is generated.
This makes representation design central rather than incidental.

I also learned that bridging between proof-friendly and computation-friendly representations is
expensive.
The \texttt{DFinsupp} layer is much nicer for mathematical statements than raw arrays, but arrays
are much easier to evaluate.
The final hybrid system works because each layer is used for what it is good at, but the conversion
and certification bridge is real overhead.
This overhead is not just programming work; it shapes the trusted statement itself.

Another lesson is that proof-producing computation often wants a checker, not a proof-producing
algorithm.
It is tempting to ask Buchberger's algorithm to produce a proof of its own correctness as it runs.
In the end, a simpler architecture was more successful: let an untrusted procedure search for the
answer, then check a finite certificate.
This is less romantic, but much more robust.

Finally, large language models were useful, but not magical.
They helped explore designs, synthesize routine proof attempts, and keep the project moving through
tedious refactors.
However, the important decisions still required mathematical and engineering judgment.
In particular, recognizing when to abandon a beautiful but non-executable proof path was a human
decision.

\section{Future Work}

\subsection{Immediate Extensions}

The most natural next step is a polished ideal-membership checker.
Given a certified \Groebner{} basis, it should be possible to reduce a target polynomial and, when
the remainder is zero, produce a Lean proof that the target lies in the ideal.
This would be a satisfying capstone because ideal membership is one of the standard first uses of
\Groebner{} bases, and the certificate machinery already points in this direction.

The user interface could also be improved substantially.
At present, examples are written directly against the project-specific polynomial types.
A better interface would include prettier printing, more convenient input syntax, and eventually a
tactic-like command that computes a basis while leaving the certificate visible in the Lean file.
Stronger negative tests would also be useful, especially examples where low fuel or malformed
witness data causes the certificate check to fail.

\subsection{Better Integration with Mathlib}

A longer-term goal is to connect the project more tightly to Mathlib.
This would involve transporting results from \texttt{CMvPolynomial} back to Mathlib's
\texttt{MvPolynomial}, relating the project's certificate predicate to standard ideal-theoretic
statements, and reconciling the project's monomial-order infrastructure with Mathlib's
\texttt{MonomialOrder}.
Doing this well would make the project much more useful for downstream formalization.

There is also more work to do on executable monomial orders.
The current tag system is a practical compromise, but it would be better to have a general pattern
for writing definitions that are both generic over monomial orders and reliably executable after
specialization.
This is especially important if one wants to add weighted orders or more elaborate order
constructions.

\subsection{External Computer Algebra Certificates}

The most practical long-term direction may be to use an external computer algebra system such as
Macaulay2 or Singular as the search engine.
Lean would send a polynomial problem to the external system, receive a candidate basis and witness
data, and then check the certificate internally.
This would preserve the trust story: Lean would not trust Macaulay2's kernel, only the certificate
that Macaulay2 helped produce.

This approach is also realistic about performance.
Computer algebra systems have decades of optimized algorithms behind them, including much faster
variants of Buchberger's algorithm.
Reimplementing all of that inside Lean is probably not the best use of effort.
A checked interface to external computation would let Lean benefit from existing software while
retaining formal guarantees for the results that matter.

\section{Conclusion}

Fettuccine began as an ambitious attempt to build verified algorithms in algebraic geometry.
It did not reach the full vision of the original proposal.
There are no Betti tables, no free resolutions, and no polished tactics for high-level algebraic
goals.
What it did produce is a coherent prototype for proof-producing \Groebner{} basis computation.

The central contribution is the architecture.
Fast, untrusted computation proposes a result; Lean checks a small certificate in a more
mathematical layer; and the final object contains an ordinary proof field.
This is enough to demonstrate that computer-algebra-style workflows can be made compatible with
machine-checked proof, even if doing so inside Lean requires careful compromises.
The project also clarified the main obstacle for future work: not the algebra alone, but the
interaction between algebraic abstraction, executable representation, and proof-producing
computation.

\bibliographystyle{ACM-Reference-Format}
\bibliography{base}

\end{document}
